\subsection*{求导为何变成乘法}
从 $\int_0^\infty y'(t)\mathrm e^{-st}\,dt$ 出发，对乘积做分部积分。导数从 $y$ 移到指数核时带出 $s$，下端边界留下 $-y(0)$，无穷远边界在适当增长与收敛条件下消失。于是得到 $\mathcal L(y')=sY-y(0)$。

再做一次，就得到二阶导数的变换，除 $s^2Y$ 外，还必须保留初始位置和初始速度。初值并没有因换到变换域而消失，而是以边界项的形式编码进去。

\subsection*{把一条完整求解链走完}
考虑 $y'+2y=0,y(0)=3$。变换后是 $sY-3+2Y=0$，解得 $Y=3/(s+2)$，再从已知指数变换识别 $y=3\mathrm e^{-2t}$。最后求导代回并检查零时刻值，完成闭环验证。

二阶受迫振子类似：先整理成 $(s^2+\omega^2)Y=F+sy_0+v_0$，再除以分母。每一项都有来源，若漏了初始值，得到的就会是另一个问题。分母中的根提示系统模态，但还需结合收敛域和约定理解。

变换法把困难移到了另一个表示中，并没有免除存在性和反演条件。Lean 在这里核验的是已经得到代数方程之后的求商步骤；积分分部与函数空间条件仍属于正文分析。读证明时要知道每一步解决了哪一部分，而不是把整个方法看成一条自动规则。
